% Pursuing Stacks (cotes 134-2 to 134-8): archive cote 134-3, site page 65 (PDF page 66), typescript page 65.
% typescript. Draft compared: ps2.tex:3353-3425.
% AI-transcribed by Claude Opus 5.5 (claude-opus-5-5, adaptive thinking, effort high), Message Batches API, batch ps1 (msgbatch_01Rv5hCaNTtUt4JHx1yqoHPL), 2026-11-12; not yet checked by eye.
\scanpage{134-3/65}
\mspage{65}

a)\quad The category A is a test c ategory in the wide sense (cf cor.1)

b)\quad Both functors $i^*$ : (Cat)\hand{$\longrightarrow$} A$^\wedge$ and $i_A$ : A$^\wedge$\hand{$\longrightarrow$}(Cat) are "modeli-%
zing", namely the set of weak equivalences in the source category
is the inverse image of the corresponding set of arrows in the
target category, and the functor induced on the localizations
with respect to weak equivalences is an equivalence of categories.

c)\quad Let \inserted[typed]{$W$ (resp.}\hand{(}$W_A$\hand{)}be the set of weak equivalences \inserted[typed]{in(Cat) (resp. \xout{\unsure{in}}}in A$^\wedge$). Then the functor
\[ (\mathrm{Hot}) \overset{\mathrm{def}}{=} W^{-1}(\mathrm{Cat}) \hand{$\longrightarrow$} W_A{}^{-1}A^\wedge \]
induced by $i^*$ is canonically isomorphic to the quasi-inverse
of the functor in opposite direction induced by $i_A$ , or equivalent%
ly, \xout{the}\inserted{can.} isomorphic to the functor in the same direction in-%
duced by $i_o^* = j_A$ : (Cat)\hand{$\longrightarrow$} A$^\wedge$(cf. again cor. 1 above for \xout{notation}
the notations).

Of course a) follows from cor. 1, and implies that $i_A$ has the
properties stated in b). That the analoguous properties hold for $i^*$ too,
and the rest of the statement i.e. c), follows formally, using the fact
that the compositum $i_Ai^*$ is canonically isomorphic to the identity functor
once we pass to the localized category (Hot) (using the theorem,
b (iii) \hand{$\Rightarrow$} a)\hand{)}.

This corollary shows that, up to canonical isomorphism, the
functor
\[ \hand{(}\mathrm{Hot}) \hand{$\longrightarrow$} W_A{}^{-1}A^\wedge \quad \text{induced by} \quad i^* : (\mathrm{Cat}) \hand{$\longrightarrow$} A^\wedge \]
\uline{does not depend on the choice of the functor} i: A \hand{$\to$}(Cat), provided
only this functor satisfies the two conditions that \xout{the} it takes it's
values in the \inserted{full} subcategory of \hand{(}Cat) of all categories with final object,
and that moreover $i^*$(\hand{$\Delta$}$_1$) be aspheric over the final object of A$^\wedge$ (plus
of course the condition of asphericity on A itself). There \uline{is} of course
always \inserted{the} \deleted{a} \uline{canonical} choice of a functor i : A \hand{$\to$}(Cat), namely $i_A$: a\hand{$\mapsto$}A$_{/a}$\unsure{,}
which (because \inserted{of} it's canonicity) looks as the best choice theoretically -
and it was the first one indeed we investigated into. But in practical
terms, the categories A$_{/a}$ are \hand{(}in \xout{most} the concrete cases one might
think of) comparatively big \xout{and} (for instance, infinite) and the corres-%
ponding functor $j_A = i_A^*$ \inserted[typed]{gives} comparatively clumsy "models" in A$^\wedge$\hand{/}for describing
the homotopy types\xout{xx} of given \xout{categories}"models" in (Cat), whereas we
can get away with considerably neater models in A$^\wedge$, using a functor
i giving rise to categories i(a) which are a lot easier to compute with,
(\deleted{and} for instance\hand{,} finite categories of very specific type). The most com\xout{\unsure{mx}}-%
monly used is of course the \uline{nerve functor} $i^*$ , corresponding to the stan-%
dar\over{t}{d} imbedding of A =\hand{$\Delta$} into (Cat) - and in the general case of the
theorem above, complemented by cor.2, \uline{the functor} $i^*$ \uline{should be viewed as}
\uline{a generalized} \uline{nerve functor}.
